\section{目标函数与噪声模型的对应：从生成机制到优化目标}

前面两篇区分了损失、风险、目标与指标，但留下一个更根本的问题：为什么回归用 MSE、分类用交叉熵？这些选择不是任意的，而是由数据的噪声结构、训练目标和正则化约束共同决定的。目标函数是``对世界的假设''的数学形式——似然项对应噪声模型，正则项对应先验约束，两者合起来才是完整的优化目标。

\subsection{先感受困难：一个异常值毁掉整个回归}

考虑简单线性回归 $y_i=\beta x_i+\varepsilon_i$，用 MSE 估计 $\beta$。若数据干净，MSE 给出最优估计；但若混入一个异常值 $(x_{n+1},y_{n+1})=(100,1000)$，MSE 会把它当作``信息量最大的点''——残差平方 $(y-\beta x)^2$ 随距离平方增长，异常值主导目标函数，把 $\widehat\beta$ 拉向自己。

MAE（绝对损失）对这个异常值不敏感：残差线性增长，异常值的影响与它的距离成正比，不会被平方放大。但 MAE 在干净数据上效率低于 MSE——对 Gaussian 噪声，MSE 是最优的（Cramér--Rao 下界），MAE 只是鲁棒的折中。

这个例子揭示了目标函数选择的根本困境：没有``普适最优''的损失，只有``与噪声模型匹配''的损失。MSE 对 Gaussian 噪声最优，MAE 对 Laplace 噪声最优，Huber 对 Gaussian+异常值混合最优。选择损失函数，就是选择对噪声的假设。

\subsection{似然项：噪声模型决定目标函数的核心}

噪声模型决定目标函数的机制是最大似然。设 $Y\mid X=x$ 的条件分布为 $p(y\mid x;\theta)$，对独立样本 $(x_i,y_i)$，似然为

\[
L(\theta)=\prod_{i=1}^n p(y_i\mid x_i;\theta).
\]

最大化似然等价于最小化负对数似然

\[
-\log L(\theta)=-\sum_{i=1}^n\log p(y_i\mid x_i;\theta).
\]

不同的噪声模型给出不同的负对数似然：

\begin{itemize}
\tightlist
\item
  \textbf{Gaussian 噪声}：$Y=\mu(x)+\varepsilon$，$\varepsilon\sim\mathcal N(0,\sigma^2)$，负对数似然是 $\frac{1}{2\sigma^2}\sum(y_i-\mu(x_i))^2+\text{常数}$——这正是 MSE（差一个常数因子）；
\item
  \textbf{Bernoulli 噪声}：$Y\sim\operatorname{Bernoulli}(p(x))$，负对数似然是 $-\sum[y_i\log p(x_i)+(1-y_i)\log(1-p(x_i))]$——这正是交叉熵；
\item
  \textbf{Laplace 噪声}：$\varepsilon\sim\operatorname{Laplace}(0,b)$，负对数似然是 $\frac1b\sum|y_i-\mu(x_i)|+\text{常数}$——这正是 MAE；
\item
  \textbf{Poisson 噪声}：$Y\sim\operatorname{Poisson}(\lambda(x))$，负对数似然是 $\sum[\lambda(x_i)-y_i\log\lambda(x_i)]$——这是 Poisson 损失。
\end{itemize}

似然项不是``评价标准''，而是``对世界的假设''的数学形式。选择 MSE 就是假设 Gaussian 噪声，选择交叉熵就是假设 Bernoulli 噪声。这个对应不是查表，而是变分推断的必然结果：给定生成模型，最大似然给出似然项。

\subsection{训练目标：似然项只是起点}

似然项由噪声模型决定，但训练目标可以更进一步：预测均值、中位数、分位数、概率还是排序，由任务需求决定。

\begin{itemize}
\tightlist
\item
  \textbf{预测均值}：MSE 对应条件均值 $\E[Y|X]$，Gaussian 噪声下最优；
\item
  \textbf{预测中位数}：MAE 对应条件中位数，Laplace 噪声下最优，对异常值鲁棒；
\item
  \textbf{预测分位数}：分位数损失（pinball loss）对应条件分位数，不对称代价（如缺货 vs 报废）下的最优选择；
\item
  \textbf{预测概率}：交叉熵对应条件概率 $P(Y|X)$，Bernoulli 噪声下最优，proper scoring rule 保证校准；
\item
  \textbf{预测排序}：pairwise loss（如 RankNet）对应相对顺序，不关心绝对值，只关心``谁排在谁前面''。
\end{itemize}

训练目标超越噪声模型：同样的 Gaussian 噪声，可以选择预测均值（MSE）、中位数（MAE）或分位数（pinball loss）。似然项给出``与噪声匹配''的目标，训练目标给出``与任务匹配''的目标。

\subsection{正则化项：先验约束优化结果}

似然项拟合数据，正则化项约束模型。完整的目标函数是

\[
J(\theta)=-\log L(\theta)+\lambda\Omega(\theta),
\]

其中 $\Omega(\theta)$ 是正则化项，$\lambda$ 控制强度。正则化不只是``防止过拟合''，而是把领域知识编码进优化：

\begin{itemize}
\tightlist
\item
  \textbf{L2（Ridge）}：$\Omega(\theta)=\norm\theta_2^2$，假设权重小（平滑），改善病态方向，对应 Gaussian 先验；
\item
  \textbf{L1（Lasso）}：$\Omega(\theta)=\norm\theta_1$，假设权重稀疏（特征选择），产生零权重，对应 Laplace 先验；
\item
  \textbf{Elastic Net}：$\Omega(\theta)=\alpha\norm\theta_1+(1-\alpha)\norm\theta_2^2$，L1 与 L2 折中，兼顾稀疏与平滑；
\item
  \textbf{早停}：训练到验证误差最小就停止，隐式正则化，防止过度拟合噪声；
\item
  \textbf{Dropout}：训练时随机失活神经元，测试时平均，假设特征不应共适应（鲁棒性）。
\end{itemize}

MAP（最大后验）视角统一了似然与正则化：

\[
\theta_{\text{MAP}}=\arg\max_\theta\log p(y|x;\theta)+\log p(\theta),
\]

其中 $\log p(\theta)$ 是先验。Gaussian 先验对应 L2，Laplace 先验对应 L1。正则化项是``对模型的假设''，与似然项的``对噪声的假设''并列。

\subsection{三个维度的统一：噪声、任务、约束}

目标函数的选择是三个维度的平衡：

\begin{enumerate}
\def\labelenumi{\arabic{enumi}.}
\tightlist
\item
  \textbf{噪声模型}：似然项匹配数据生成机制（Gaussian→MSE，Bernoulli→交叉熵，Laplace→MAE）；
\item
  \textbf{训练目标}：似然项之外，任务需求决定预测什么（均值、中位数、概率、排序）；
\item
  \textbf{正则化约束}：先验或领域知识约束模型（L2 平滑、L1 稀疏、早停、Dropout）。
\end{enumerate}

三个维度可以独立变化：Gaussian 噪声下可以选择预测均值（MSE）或中位数（MAE），可以加 L2 或 L1 正则化。但维度之间也有交互：强正则化可能偏离似然项的最优解（偏差-方差权衡），任务目标可能要求特定的似然形式（分位数预测需要分位数损失，不能用 MSE 替代）。

\subsection{目标函数的形状决定优化几何}

目标函数的形状（凸性、光滑性、增长速率）决定优化的难度与性质：

\begin{itemize}
\tightlist
\item
  \textbf{MSE 是二次的}：梯度线性，Hessian 常数，凸优化，线性回归有闭式解（正规方程）。但二次增长对异常值敏感——残差平方随距离平方增长；
\item
  \textbf{交叉熵是凸的}：梯度含 sigmoid/softmax，无闭式解，需要迭代优化（梯度下降、Newton）。凸性保证全局最优，但条件数可能很差（类别不平衡时）；
\item
  \textbf{MAE 是分段线性的}：次梯度（在零点不可导），中位数回归，对异常值鲁棒。但线性增长意味着对 Gaussian 噪声效率低（渐近方差比 MSE 大 $\pi/2$ 倍）；
\item
  \textbf{Hinge 是分段线性的}：次梯度，支持向量（稀疏解），几何间隔最大化。不可导点对应决策边界，次梯度提供下降方向。
\end{itemize}

形状决定优化几何：二次损失给出光滑凸问题，分段线性损失给出非光滑凸问题，0-1 损失给出非凸组合问题。选择损失函数，就是选择优化的几何。

\subsection{目标函数的统计性质决定泛化}

目标函数的统计性质（偏差-方差、校准、proper scoring rule）决定泛化能力：

\begin{itemize}
\tightlist
\item
  \textbf{偏差-方差}：MSE 对 Gaussian 噪声是有效的（达到 Cramér--Rao 下界），但对异常值高方差；MAE 对异常值低方差，但对 Gaussian 噪声低效。偏差-方差权衡在目标函数选择中已经体现；
\item
  \textbf{校准}：交叉熵是 strictly proper scoring rule——最优预测是真实条件概率 $p(y|x)$，不是``最能分开两类的分数''。这保证概率输出可信，可以直接用于决策阈值（Part 11 的统计决策）；
\item
  \textbf{替代损失}：0-1 损失是最优的分类目标（Bayes 风险），但不可导。Hinge、logistic、交叉熵都是凸替代——在一定条件下（如 $p(y|x)$ 单调），替代损失的最优解与 0-1 损失一致，但替代损失本身可能过度惩罚``容易分类的点''（logistic）或``边界附近的点''（hinge）。
\end{itemize}

统计性质决定泛化：MSE 的平方增长给出低偏差但高方差，MAE 的线性增长给出高偏差但低方差，交叉熵的 proper scoring 给出校准概率。

\subsection{边界清单}

\begin{itemize}
\tightlist
\item
  目标函数与噪声模型匹配时最优，错配时失效：MSE 对 Gaussian 噪声最优，但对 Laplace 噪声低效；交叉熵对 Bernoulli 噪声最优，但对重尾噪声（如 t 分布）可能过度自信；
\item
  模型错设时目标函数失效：若真实噪声是混合分布，任何单一噪声假设的目标函数都只是近似；鲁棒方法（Huber、trimmed loss）是折中，不是最优；
\item
  正则化强度 $\lambda$ 是超参数，需要交叉验证选择：过强的正则化导致欠拟合，过弱的正则化导致过拟合；
\item
  替代损失与真实损失的差距：0-1 损失不可导，凸替代（hinge、logistic）在一定条件下一致，但替代损失本身可能过度惩罚某些点；
\item
  目标函数与优化算法交互：MSE+线性模型=闭式解，MSE+深度网络=非凸优化；交叉熵+线性模型=凸优化，交叉熵+深度网络=非凸优化。
\end{itemize}

\subsection{落点}

目标函数的选择是三个维度的统一：似然项匹配噪声模型（Gaussian→MSE，Bernoulli→交叉熵），训练目标超越噪声（预测均值、中位数、概率、排序），正则化项约束模型（L2 平滑、L1 稀疏、早停、Dropout）。这不是查表，而是最大似然与最大后验的必然结果——噪声模型给出似然，先验给出正则化，任务需求选择预测目标。目标函数的形状决定优化几何，统计性质决定泛化，选择目标函数就是选择任务定义。这是从``会用损失函数''到``理解损失函数''的关键一步。
